% ===== PRÉAMBULE CANONIQUE — à coller VERBATIM en tête de CHAQUE .tex =====
\documentclass[11pt,a4paper]{article}
\usepackage[utf8]{inputenc}\usepackage[T1]{fontenc}\usepackage{lmodern}
\usepackage[french]{babel}
\usepackage{amsmath,amssymb,mathtools}\usepackage{icomma}
\usepackage[version=4]{mhchem}          % \ce{...} : équations & formules
\usepackage{siunitx}\sisetup{locale=FR} % \qty{}{}, \si{}, \ang{}
\usepackage{chemfig}                    % \chemfig{...} : structures développées
\usepackage{xcolor}\usepackage{colortbl}
\usepackage{tikz}\usepackage{pgfplots}\pgfplotsset{compat=1.18}
\usepackage[most]{tcolorbox}\usepackage{enumitem}\usepackage{array,booktabs}
\usepackage[a4paper,margin=2.2cm]{geometry}
\usetikzlibrary{arrows.meta,calc,angles,quotes,positioning,patterns,backgrounds,shapes.geometric}
\definecolor{primary}{RGB}{222,49,99}\definecolor{primarydark}{RGB}{150,22,55}
\definecolor{defcol}{RGB}{0,114,178}\definecolor{methcol}{RGB}{52,92,148}
\definecolor{excol}{RGB}{0,135,90}\definecolor{attncol}{RGB}{200,40,55}
\tcbset{boxbase/.style={enhanced,breakable,boxrule=0.9pt,arc=2.5pt,
  left=3mm,right=3mm,top=2.3mm,bottom=2.3mm,fonttitle=\bfseries,coltitle=white}}
% --- boîtes TUTORIEL (noms EXACTS, ne pas renommer) ---
\newtcolorbox{cle}[1][]{boxbase,colback=defcol!6,colframe=defcol,
  title={\textsc{À retenir}\ifx\relax#1\relax\relax\else\textmd{ : \itshape #1}\fi}}
\newtcolorbox{methode}[1][]{boxbase,colback=methcol!7,colframe=methcol,sharp corners,
  title={\textsc{Méthode}\ifx\relax#1\relax\relax\else\textmd{ --- \itshape #1}\fi}}
\newtcolorbox{exemplebox}[1][]{boxbase,colback=excol!5,colframe=excol,
  title={\textsc{Exemple}\ifx\relax#1\relax\relax\else\textmd{ : \itshape #1}\fi}}
\newtcolorbox{piege}[1][]{boxbase,colback=attncol!6,colframe=attncol,
  title={\textsc{Piège}\ifx\relax#1\relax\relax\else\textmd{ : \itshape #1}\fi}}
\newtcolorbox{remarque}[1][]{boxbase,colback=black!4,colframe=black!45,
  title={\textsc{Remarque}\ifx\relax#1\relax\relax\else\textmd{ : \itshape #1}\fi}}
% --- boîtes EXERCICE / CORRIGÉ (noms EXACTS, auto-comptées) ---
\newtcolorbox[auto counter]{exo}[1][]{boxbase,colback=primary!4,colframe=primary,
  title={\textsc{Exercice}~\thetcbcounter\ifx\relax#1\relax\relax\else\textmd{ --- \itshape #1}\fi}}
\newtcolorbox[auto counter]{corrige}[1][]{boxbase,colback=excol!4,colframe=excol,
  title={\textsc{Corrigé}~\thetcbcounter\ifx\relax#1\relax\relax\else\textmd{ --- \itshape #1}\fi}}
\newcommand{\R}{\mathbb{R}}\newcommand{\N}{\mathbb{N}}\newcommand{\Z}{\mathbb{Z}}\newcommand{\ds}{\displaystyle}
% (un \usepackage{fancyhdr}+en-tête est possible ; il est ignoré côté web)
% =========================================================================
\begin{document}
\sisetup{per-mode=symbol}
\begin{corrige}[le compte-rendu du laboratoire]
Une mesure a un nombre \emph{fini} de CS et son dernier chiffre écrit est l'incertain ; un nombre
compté est exact (infinité de CS, aucun chiffre incertain).
\begin{enumerate}[label=\alph*)]
  \item $\qty{18,60}{\milli\litre}$ : $\mathbf{4}$ CS ; chiffre incertain $=$ le $0$ final.
  \item $\qty{0,04060}{\gram}$ : zéros de tête non sig.\ ; le $0$ captif et le $0$ de queue comptent
        $\Rightarrow \mathbf{4}$ CS ; chiffre incertain $=$ le $0$ final.
  \item $3$ pastilles \emph{comptées} : nombre \textbf{exact}, infinité de CS, \emph{pas} de chiffre
        incertain.
  \item $\num{5,0e-3}\ \si{\mol\per\litre}$ : $\mathbf{2}$ CS ; chiffre incertain $=$ le $0$ de $5,0$.
\end{enumerate}
\end{corrige}

\begin{corrige}[lever une ambiguïté en contexte]
\begin{enumerate}[label=\alph*)]
  \item «~$\qty{25}{\kilo\gram}$~» est une indication ambiguë (au mieux $2$ CS) : le zéro final d'un
        entier n'est pas fiable. «~$\qty{25,00}{\kilo\gram}$~» affiche $\mathbf{4}$ CS. La balance nous
        apprend que la masse est en réalité connue \emph{au centième de kilogramme} ($\pm\qty{0,01}{\kilo\gram}$),
        soit une précision bien supérieure.
  \item «~au dixième de kilogramme près~» signifie que le dernier chiffre fiable est celui des dixièmes :
        on écrit $\qty{25,0}{\kilo\gram}$ ($3$ CS), soit en notation scientifique
        $\boxed{2,50\times10^{1}\ \si{\kilo\gram}}$.
\end{enumerate}
\end{corrige}

\begin{corrige}[vrai ou faux, justifié]
\begin{enumerate}[label=\alph*)]
  \item \textbf{Faux.} $0,0025$ a $2$ CS : les zéros de tête ne comptent pas.
  \item \textbf{Faux.} $\qty{4,50}{\gram}$ et $\qty{0,00450}{\kilo\gram}$ ont \emph{tous deux} $3$ CS :
        changer d'unité ne change jamais le nombre de CS.
  \item \textbf{Vrai.} Un nombre exact (comptage ou définition) possède une infinité de CS.
  \item \textbf{Faux.} $\num{3,0e2}$ a $2$ CS, tandis que $300$ est \emph{ambigu} ($1$, $2$ ou $3$ CS) :
        ils n'ont pas forcément le même nombre de CS. C'est justement pourquoi on préfère la notation
        scientifique.
\end{enumerate}
\end{corrige}

\begin{corrige}[le zéro qui trahit]
\begin{enumerate}[label=\alph*)]
  \item $\qty{7,0}{\gram}$ : $\mathbf{2}$ CS ; $\qty{7}{\gram}$ : $\mathbf{1}$ CS.
  \item L'étudiant a perdu l'information «~mesure fiable au dixième de gramme~» : le zéro de queue était
        \emph{porteur d'information}. «~$\qty{7}{\gram}$~» ne promet plus qu'une précision à $\pm\qty{1}{\gram}$.
  \item Mesure correcte : $\boxed{7,0\times10^{0}\ \si{\gram}}$ (soit $\qty{7,0}{\gram}$, $2$ CS).
\end{enumerate}
\end{corrige}

\end{document}
